\documentclass[10pt,a4paper]{amsart}
\usepackage[margin=19mm]{geometry}
\usepackage{amsmath,amssymb,mathtools}
\usepackage[scr=rsfs,cal=euler]{mathalfa}
\usepackage{xfrac}
\usepackage[unicode,hidelinks,bookmarks=false]{hyperref}
\newcommand{\ie}{i.e.\ }
\newcommand{\cf}{cf.\ }
\newcommand{\resp}{respectively\ }
\newcommand{\Subsection}{\subsection}
\providecommand{\nobreakdash}{\nobreak\hbox{-}}
\setlength{\emergencystretch}{2em}
\multlinegap0pt
\theoremstyle{plain}
\newtheorem{theorem}{Theorem}
\newtheorem{lemma}{Lemma}
\newtheorem{conjecture}{Conjecture}
\newtheorem{proposition}{Proposition}
\newtheorem{claim}{Claim}
\theoremstyle{definition}
\newtheorem{definition}{Definition}
\theoremstyle{remark}
\newtheorem{remark}{Remark}
\newcommand{\minorOf}{\leq_{M}}
\newcommand{\IN}{\mathbb{N}}
\newcommand{\therese}[1]{}\newcommand{\dinis}[1]{}
\title[Bipartite Graphs Are Not Well-Quasi-Ordered by Bipartite Minors]{Bipartite Graphs Are Not Well-Quasi-Ordered by Bipartite Minors: the tadpole is a bipartite minor but not a minor of the cycle (printed as Theorem 1)}
\author{Therese Biedl, Dinis Vitorino}
\date{Source: arXiv:2601.23101v3 (CC BY 4.0) (CC BY 4.0)}
\begin{document}
\maketitle

\noindent\emph{Source and licence.} Bipartite Graphs Are Not Well-Quasi-Ordered by Bipartite Minors. Version arXiv:2601.23101v3 (CC BY 4.0), distributed under the Creative Commons Attribution 4.0 International licence, http://creativecommons.org/licenses/by/4.0/, as recorded in the arXiv licence metadata for this exact version and shown on the saved abstract page https://arxiv.org/abs/2601.23101v3. Full licence notice: licenses/arxiv-2601.23101-CC-BY-4.0.txt.\par\smallskip

\noindent\textit{Editorial scope.} Counted unit: the tadpole theorem, printed in the article as Theorem 1 (source0.tex lines 271--274, one \texttt{theorem} environment with source label \texttt{thm: bip minor but not minor}), delivered together with the article's complete original proof of it (source0.tex lines 276--300, one \texttt{proof} environment of 25 lines immediately adjacent to the statement, with no gap and no deferral). The proof is the article's own argument: it exhibits the explicit sequence of admissible contractions that turns the cycle into the tadpole with the required head and pole, and then rules out the reverse relation by observing that the cycle has maximum degree two, that a graph of maximum degree three is the tadpole's target, and that the graphs of maximum degree at most two form a minor-closed class. No mathematics is written, repaired or re-derived: statement and proof are the article's own text line for line, in the article's own notation, and no line of either is omitted, substituted or re-flowed.

Required context retained uncounted: the article's own sentence defining the tadpole and fixing the notation $T(l,l_1)$ (source0.tex line 246). Nothing else from the article is needed.

References. No citation command occurs in any retained line, so this package carries no bibliography block and no bibliography entry.

Editorial formatting only, in addition: an AMS \texttt{amsart} document that restates verbatim the article's own macro definitions needed by the retained lines, including the article's own two working-comment macros in the form in which they are effective in the article (the article switches them to print nothing, so the brief comments the authors address to each other inside the proof print nothing here either), and the article's own theorem-environment declarations. The article's own publisher class file \texttt{elsarticle.cls}, its bibliography style file and its bibliography database are neither shipped in this package nor used to build it. Consequently the retained environments print here with independent counters, so the counted theorem prints as Theorem 1, which is also how the article prints it (the article's numbering of its three theorems is preserved because the delivered theorem is the article's first); the equations of the retained spans are numbered anew in order of appearance and every cross-reference inside the retained text resolves to this wrapper's numbering. That is the only change to the numbering. No figure, no image-inclusion line and no drawing code occur in any retained line, so no artwork is delivered or needed, although the article contains figures elsewhere.

Not retained: the article's abstract, its introduction, its definitions and figures, its other two theorems and its open problems; none of that material is used as a step of the delivered proof, and the delivered proof is complete as the author wrote it. The package ships the article's concatenated source verbatim as \texttt{sources/193683/source0.tex}, and every delivered excerpt is an exact line range of that file. No mathematical statement, hypothesis, estimate or conclusion has been rewritten, repaired or added, and printed slips, if any, are preserved literally.
The \emph{tadpole} $T$ with a \emph{head} of length $l$ and a \emph{pole} of length $l_1$ is denoted by $T(l,l_1)$, and is the graph obtained from a cycle $C_l$ and a path $P_{l_1}$ by picking a vertex in $C_l$ and adding an edge between it and an endpoint of $P_{l_1}$.

\begin{theorem}
    \label{thm: bip minor but not minor}
    For all integers $l\geq 3$ and $l_1\geq 1$, the tadpole $T(l,l_1)$ is a bipartite minor of $C_{p}$ for any $p\geq l+2l_1$ with the same parity as $l$, but $T(l,l_1)$ is not a minor of $C_p$ for any integer $p\geq 3$.
\end{theorem}

\begin{proof}
    Consider the graph $C_p$ where $p=l+2k$ and $k\geq l_1$.
    Contract a pair of vertices of $C_p$ that have a common neighbour;
    this yields $T(p-2,1)$.
    Let $w$ be the vertex of the head of $T(p-2,1)$ with a neighbour on its pole, and let $u,v$ be its two neighbours on the head.
    The contraction of $u,v$ is admissible and yields $T(p-4,2)$.
    \therese{switched $v$ and $w$ here to match how we defined contractions earlier}
    Repeating $k-2$ times (always contracting the two vertices of the head that have the vertex adjacent to the pole as a common neighbour) gives $T(p-2k,k)=T(l,k)$.
    Deleting the $k-l_1$ vertices at the end of the pole that is furthest from the head now gives $T(l,l_1)$ as desired.

    % To see that $T(l,l_1)\not\minorOf C_{p}$ for any $p\in\IN_{\geq 3}$, note that all vertices $C_p$ have degree at most two, and it is easy to see that the same holds for all its minors.
    % %\footnote{Although contractions do not generally preserve maximum degree, this is true when the maximum degree is two.}.
    % \therese{I think this can be safely stated as obvious}
    % \dinis{I think the main reason I changed this was that the reviewer thought the original wording might lead one to believe that all minors of graphs with maximum degree $\Delta$ have maximum degree $\Delta$.
    % Fearing that this problem persists, I added a footnote (commented out now), but maybe that's unnecessary.}
    % As $T(l,l_1)$ has a vertex of degree three, it cannot be a minor of $C_{p}$.
    To see that $T(l,l_1)\not\minorOf C_{p}$ for any $p\in\IN_{\geq 3}$, one need only observe that $C_p$ has maximum degree two, but $T(l,l_1)$ has a vertex of degree three.
    As one can easily check, the graphs with maximum degree at most two form a minor-closed class (although taking minors does not generally preserve maximum degree), yielding the desired claim.
    \dinis{I reworded a bit here (see discussion/commented out things above)}
    \therese{I don't see any rewording except the `subquadratic', but I think what's here is fine.}
    \therese{I had not heard the term `subquadratic graph' before, and would have guessed that this means `has asymptotically fewer than $n^2$ edges'.  Let's not use this term.  (Or if you really love it, define it.)}
    \dinis{Fair.
    I'm not sure the term subquadratic exists anywhere else.
    I mostly wanted to avoid using the words maximum degree and minor-closed to prevent people from thinking that any upper bound on max degree is preserved by taking minors (which is obviously false)}
\end{proof}

\end{document}
